% TOP_PROBLEM: {"id": 20001054, "title": "Finite-boundary transfer for infinite second neighborhoods", "category_id": 2, "category": {"id": 2, "name": "combinatorics", "display_name": "Combinatorics", "description": "Counting problems, graph theory, discrete structures.", "slug": "combinatorics", "order_index": 2, "created_at": "2026-07-31T15:26:25.671Z"}, "status": "partially_solved", "problem_number": "AIM-COMBINATORICS-0179", "set_id": 14}
% FIRST_POSTED: 2026-10-01
% ENTRY_KIND: statement_only
% UnsolvedMath ID 20001054; source catalogue code AIM-COMBINATORICS-0179
% !TeX program = lualatex
% Definition and sources only; this file is not an LLM solution attempt.
\documentclass[11pt,a4paper]{article}
\usepackage[margin=25mm]{geometry}
\usepackage{fontspec}
\setmainfont{lmroman10-regular.otf}[BoldFont=lmroman10-bold.otf,
 ItalicFont=lmroman10-italic.otf,BoldItalicFont=lmroman10-bolditalic.otf]
\usepackage{amsmath,amssymb,mathtools}
\usepackage{unicode-math}
\setmathfont{latinmodern-math.otf}
\AtBeginDocument{\let\setminus\smallsetminus}
\usepackage{xurl,hyperref}
\hypersetup{colorlinks=true,urlcolor=blue}
\setcounter{secnumdepth}{0}
\setlength{\parindent}{0pt}
\setlength{\parskip}{5pt}
\setlength{\emergencystretch}{3em}
\begin{document}
\section*{Finite-boundary transfer for infinite second neighborhoods}\label{20001054}
{\small UnsolvedMath ID 20001054 · AIM-COMBINATORICS-0179 · Combinatorics · AIM Workshop Problem Lists}
\subsection{Definitions and mathematical statement}
Let $D=(V,E)$ be a nonempty, possibly infinite, oriented graph: there
are no loops or parallel arcs and never arcs in both directions
between two vertices. For a vertex $v$, let $N^+(v)$ and $N^-(v)$
be its sets of outgoing and incoming neighbours. Define
\[
 N_2^+(v)=\{w\notin N^+(v)\cup\{v\}:\exists u, v\to u\to w\}.
\]
Paths have finite length, and this set is precisely the vertices at
shortest outgoing distance two.
Call $D$ locally finite when both $N^+(v)$ and $N^-(v)$ are finite
for every $v$. Call it out-locally finite when only $N^+(v)$ is
required to be finite for every $v$; there need not be a uniform
bound on these finite degrees.

The workshop asks two questions. Must every nonempty locally finite
oriented graph have a vertex $v$ with
$|N_2^+(v)|\ge|N^+(v)|$? Does the same conclusion hold for every
nonempty out-locally finite oriented graph?
The latter permits infinite indegrees and is the stronger assertion.
In either class the second out-neighbourhood of a fixed vertex is
finite, since it is contained in a finite union of finite first
out-neighbourhoods. Thus the comparison is between finite integers,
not infinite cardinalities. No countability, connectedness, or global
finiteness is imposed on the vertex set. These are extensions of the
finite second-neighbourhood problem to the two stated local finiteness
conditions.
\subsection{Short English statement}
Does the second-neighbourhood conjecture remain valid for infinite locally finite graphs, or even when only each outdegree is finite?
\subsection{Sources}
\begin{itemize}
\item[S1] ULAM.AI and the UnsolvedMath Contributors, supplied problems.json, record 20001054 (AIM-COMBINATORICS-0179), AIM Workshop Problem Lists. Catalogue identity and source transcription; CC BY 4.0.
\url{https://huggingface.co/datasets/ulamai/UnsolvedMath}.
\item[S2] American Institute of Mathematics, The Caccetta-Haggkvist conjecture, item 6.11. Original workshop question underlying AIM-COMBINATORICS-0179.
\url{https://aimath.org/WWN/caccetta/caccetta.pdf}.
\end{itemize}
\end{document}
